\documentclass[12pt]{book} % Change to book
\usepackage{amsmath} % for mathematical expressions
\usepackage{xcolor} % for colored text
\usepackage{graphicx} % for including images
\usepackage{url}
\usepackage{booktabs} % for better-looking tables
\usepackage{makecell}
\usepackage{indentfirst}
\usepackage{algorithm}
\usepackage{algpseudocode}
\usepackage{amssymb}
\usepackage{placeins} % For \FloatBarrier
\usepackage{listings}
% Define the language style for Python code
\lstset{
    language=Python,
    basicstyle=\ttfamily,
    keywordstyle=\color{blue},
    stringstyle=\color{red},
    commentstyle=\color{gray},
    morecomment=[l][\color{magenta}]{\#},
    frame=single,
    breaklines=true
}
\begin{document}
\chapter{Scheduling: First Things First}

\section{Introduction to Scheduling}

Scheduling is key in computer science and operations research, focusing on efficiently allocating resources over time to complete tasks or jobs. It is vital in operating systems, manufacturing, project management, and networking. This section covers the definition, importance, and applications of scheduling, as well as common scheduling problems.

\subsection{Definition and Importance}

Scheduling allocates resources like processors, machines, or personnel to tasks over time, aiming to optimize criteria such as throughput, latency, response time, or resource utilization. It is an optimization problem where resources are assigned to tasks to meet constraints and optimize an objective function.

For example, in the \textbf{single-machine scheduling problem}, a set of jobs with given processing times and deadlines must be scheduled on a single machine to minimize total completion time or maximum lateness. Given \(n\) jobs with processing times \(p_1, p_2, \ldots, p_n\) and deadlines \(d_1, d_2, \ldots, d_n\), the goal is to minimize an objective function \(f\) under certain constraints.

Effective scheduling improves resource utilization, minimizes response times, meets deadlines, and boosts system efficiency. Efficient scheduling algorithms are essential for optimizing various systems, leading to cost savings, higher productivity, and better customer satisfaction.

\subsection{Applications in Computing and Beyond}

Scheduling techniques are applied in many fields:

\begin{itemize}
    \item \textbf{Operating Systems:} Managing process or thread execution on CPUs with algorithms like round-robin, shortest job first, and priority scheduling.
    \item \textbf{Manufacturing:} Optimizing production, allocating machines and workers, and minimizing costs with job shop and flow shop scheduling.
    \item \textbf{Networking:} Managing data packet transmission and allocating network resources with packet scheduling in routers and traffic scheduling in wireless networks.
    \item \textbf{Project Management:} Planning resources, scheduling tasks, and tracking progress with methods like the critical path method (CPM) and program evaluation and review technique (PERT).
\end{itemize}

\subsection{Overview of Scheduling Problems}

Scheduling problems vary based on resource numbers, task nature, and optimization criteria:

\begin{itemize}
    \item \textbf{Single-Machine Scheduling:} Scheduling jobs on a single machine with constraints like processing times and deadlines.
    \item \textbf{Parallel-Machine Scheduling:} Scheduling jobs on multiple identical or different machines to optimize makespan or total completion time.
    \item \textbf{Flow Shop Scheduling:} Jobs pass through a sequence of machines to minimize total completion time or makespan.
    \item \textbf{Job Shop Scheduling:} Scheduling jobs requiring multiple machines with different processing times and sequence-dependent setup times.
    \item \textbf{Resource-Constrained Project Scheduling:} Scheduling tasks with precedence relations under resource constraints and project deadlines.
\end{itemize}

Many scheduling problems are NP-hard, necessitating efficient heuristic or approximation algorithms to find near-optimal solutions in reasonable time.



\section{Theoretical Foundations of Scheduling}

Scheduling is a core problem in computer science and operations research, crucial in domains like manufacturing, transportation, and computing systems. This section covers the classification of scheduling problems, optimality criteria, complexity measures, and computability.

\subsection{Classification of Scheduling Problems}

Scheduling problems are classified based on job characteristics, processing environment, and scheduling objectives. This helps in understanding different problems and developing appropriate solutions.

\textbf{Characteristics of Jobs or Tasks}

Jobs can vary by:

\begin{itemize}
    \item \textbf{Processing times:} Jobs may have deterministic processing times (known in advance) or stochastic processing times (following a probability distribution).
    \item \textbf{Release times and due dates:} Jobs have release times (when they become available) and due dates (when they must be completed). These problems are dynamic as they involve scheduling over time.
    \item \textbf{Precedence constraints:} Some jobs must be completed before others can start, common in job shop or flow shop scheduling problems in manufacturing.
\end{itemize}

\textbf{Processing Environment}

Jobs are scheduled in different environments:

\begin{itemize}
    \item \textbf{Single-machine scheduling:} One machine processes all jobs. The goal is to determine the job order to optimize criteria like makespan or total completion time.
    \item \textbf{Parallel machine scheduling:} Multiple identical machines process jobs. The objective is to assign jobs to machines to minimize criteria like total completion time or maximum lateness.
    \item \textbf{Open-shop scheduling:} Multiple machines are available, and each job must be processed on a specific sequence of machines to optimize a criterion like makespan.
\end{itemize}

\textbf{Scheduling Objectives}

Objectives vary in scheduling problems:

\begin{itemize}
    \item \textbf{Minimization of makespan:} Total time to complete all jobs.
    \item \textbf{Minimization of total completion time:} Sum of all job completion times.
    \item \textbf{Minimization of maximum lateness:} Maximum time by which job completion exceeds its due date.
\end{itemize}

These classifications help in analyzing and solving scheduling problems efficiently.

\subsection{Optimality Criteria}

Optimality criteria define the desired outcomes in scheduling problems, like minimizing makespan, total completion time, or maximum lateness.

Given \(n\) jobs and \(m\) machines, the scheduling problem can be formulated as:

\[
\text{Minimize } f(x_1, x_2, \ldots, x_n)
\]
subject to:
\[
g(x_1, x_2, \ldots, x_n) \leq b
\]
where \(x_i\) represents the assignment of job \(i\) to a machine, and \(f\) is the objective function.

\subsection{Complexity and Computability}

Scheduling problem complexity is analyzed using computational complexity theory. For instance, consider a problem with \(n\) jobs and \(m\) machines, where each job \(j\) has a processing time \(p_j\) and must be processed on one machine without preemption.

\textbf{Proposition 1:} The number of operations required to find an optimal schedule for \(n\) jobs and \(m\) machines is \(O(n \cdot m)\).

\textbf{Proof:} 
To find an optimal schedule, all possible job-to-machine assignments must be considered. Since each job can be assigned to any of the \(m\) machines, there are \(m^n\) combinations. Thus, the operations required are \(O(n \cdot m)\).

Many scheduling problems are NP-hard, meaning no polynomial-time algorithm exists unless P=NP.

\textbf{Theorem 1:} The flow shop scheduling problem is NP-hard.

\textbf{Proof:} 
Reducing the flow shop scheduling problem to the traveling salesman problem (TSP), which is NP-hard, shows the complexity. Each city in TSP corresponds to a job, and each distance corresponds to processing time on a machine. Since TSP is NP-hard, so is flow shop scheduling.

Understanding these foundations helps in tackling scheduling problems effectively in various applications.



\section{Scheduling Algorithms}

Scheduling algorithms are essential in operating systems and computer science, determining the order in which tasks or processes are executed on a computing system. Different scheduling algorithms prioritize tasks differently based on various criteria such as waiting time, execution time, priority, etc. In this section, we will explore several common scheduling algorithms, including First-Come, First-Served (FCFS), Shortest Job First (SJF), Priority Scheduling, and Round-Robin Scheduling.

\subsection{First-Come, First-Served (FCFS)}

First-Come, First-Served (FCFS) is one of the simplest scheduling algorithms, where tasks are executed in the order they arrive in the system. It follows a non-preemptive approach, meaning once a process starts executing, it continues until completion without interruption.

Let's denote $n$ as the total number of processes, $AT_i$ as the arrival time of process $i$, $BT_i$ as the burst time of process $i$, and $CT_i$ as the completion time of process $i$. The completion time for each process can be calculated using the formula:

\[
CT_i = \sum_{j=1}^{i} BT_j
\]

The average waiting time ($\text{AWT}$) for all processes can be computed as the total waiting time divided by the number of processes ($n$):

\[
\text{AWT} = \frac{\sum_{i=1}^{n} WT_i}{n}
\]

where $WT_i$ is the waiting time for process $i$, calculated as:

\[
WT_i = CT_i - AT_i - BT_i
\]

Now, let's provide the algorithm for the FCFS technique and explain it in detail:
\FloatBarrier
\begin{algorithm}[H]
\caption{First-Come, First-Served (FCFS) Scheduling}
\begin{algorithmic}[1]
\State Sort processes based on arrival time $AT_i$
\State $CT_1 \gets AT_1 + BT_1$
\For{$i \gets 2$ to $n$}
    \State $CT_i \gets CT_{i-1} + BT_i$
\EndFor
\end{algorithmic}
\end{algorithm}
\FloatBarrier
The FCFS algorithm simply iterates through the processes in the order of their arrival and calculates the completion time for each process sequentially.

\subsection{Shortest Job First (SJF)}

Shortest Job First (SJF) is a scheduling algorithm that selects the process with the shortest burst time to execute next. It can be either preemptive or non-preemptive. In the non-preemptive version, once a process starts execution, it continues until completion without interruption.

Let's denote $n$ as the total number of processes, $AT_i$ as the arrival time of process $i$, $BT_i$ as the burst time of process $i$, and $CT_i$ as the completion time of process $i$. The completion time for each process can be calculated using the formula:

\[
CT_i = CT_{i-1} + BT_i
\]

where $CT_0 = 0$. The average waiting time ($\text{AWT}$) for all processes can be computed similarly to FCFS.

Now, let's provide the algorithm for the SJF technique and explain it in detail:
\FloatBarrier
\begin{algorithm}[H]
\caption{Shortest Job First (SJF) Scheduling}
\begin{algorithmic}[1]
\State Sort processes based on burst time $BT_i$
\State $CT_1 \gets AT_1 + BT_1$
\For{$i \gets 2$ to $n$}
    \State $CT_i \gets CT_{i-1} + BT_i$
\EndFor
\end{algorithmic}
\end{algorithm}
\FloatBarrier
The SJF algorithm sorts processes based on their burst time in ascending order and then follows the same steps as FCFS to calculate completion times.

\subsection{Priority Scheduling}

Priority Scheduling is a scheduling algorithm where each process is assigned a priority, and the process with the highest priority is selected for execution first. It can be either preemptive or non-preemptive.

Let's denote $n$ as the total number of processes, $AT_i$ as the arrival time of process $i$, $BT_i$ as the burst time of process $i$, and $CT_i$ as the completion time of process $i$. The completion time for each process can be calculated similarly to FCFS.

Now, let's provide the algorithm for the Priority Scheduling technique and explain it in detail:
\FloatBarrier
\begin{algorithm}[H]
\caption{Priority Scheduling}
\begin{algorithmic}[1]
\State Assign priority to each process
\State Sort processes based on priority
\State $CT_1 \gets AT_1 + BT_1$
\For{$i \gets 2$ to $n$}
    \State $CT_i \gets CT_{i-1} + BT_i$
\EndFor
\end{algorithmic}
\end{algorithm}
\FloatBarrier
The Priority Scheduling algorithm assigns priorities to processes and then sorts them based on their priority. It then calculates completion times similarly to FCFS.

\subsection{Round-Robin Scheduling}

Round-Robin Scheduling is a preemptive scheduling algorithm where each process is assigned a fixed time quantum, and processes are executed in a circular queue. If a process does not complete within its time quantum, it is preempted and placed at the end of the queue to wait for its turn again.

Let's denote $n$ as the total number of processes, $AT_i$ as the arrival time of process $i$, $BT_i$ as the burst time of process $i$, $CT_i$ as the completion time of process $i$, and $TQ$ as the time quantum.

The completion time for each process can be calculated using the formula:

\[
CT_i = \min(CT_{i-1} + TQ, AT_i + BT_i)
\]

where $CT_0 = 0$. The average waiting time ($\text{AWT}$) for all processes can be computed similarly to FCFS.

Now, let's provide the algorithm for the Round-Robin Scheduling technique and explain it in detail:
\FloatBarrier
\begin{algorithm}[H]
\caption{Round-Robin Scheduling}
\begin{algorithmic}[1]
\State Set $CT_1 \gets \min(AT_1 + BT_1, TQ)$
\For{$i \gets 2$ to $n$}
    \State $CT_i \gets \min(CT_{i-1} + TQ, AT_i + BT_i)$
\EndFor
\end{algorithmic}
\end{algorithm}
\FloatBarrier
The Round-Robin Scheduling algorithm assigns a fixed time quantum to each process and processes them in a circular queue, calculating completion times accordingly.


\section{Advanced Scheduling Techniques}

Scheduling techniques are essential in computer science and operating systems for efficiently managing resources and executing tasks. In this section, we delve into advanced scheduling techniques, including preemptive vs non-preemptive scheduling, multiprocessor scheduling, and real-time scheduling.

\subsection{Preemptive vs Non-Preemptive Scheduling}

Preemptive and non-preemptive scheduling are two fundamental approaches to task scheduling in operating systems. Preemptive scheduling allows the scheduler to interrupt a running task to allocate CPU time to another task with higher priority, while non-preemptive scheduling does not allow such interruptions.

\subsubsection{Preemptive Scheduling Algorithm}

Preemptive scheduling algorithms prioritize tasks based on their priority levels, executing higher-priority tasks before lower-priority ones. One such algorithm is the Shortest Remaining Time First (SRTF) algorithm, which selects the task with the shortest remaining processing time whenever a new task arrives or a running task completes.
\FloatBarrier
\begin{algorithm}
\caption{Shortest Remaining Time First (SRTF)}
\begin{algorithmic}[1]
\Procedure{SRTF}{tasks}
\State $current\_time \gets 0$
\While{$tasks$ not empty}
\State $next\_task \gets$ task with shortest remaining time
\State $execute(next\_task)$
\State $current\_time \gets current\_time + next\_task.duration$
\EndWhile
\EndProcedure
\end{algorithmic}
\end{algorithm}
\FloatBarrier
\subsubsection{Non-Preemptive Scheduling Algorithm}

Non-preemptive scheduling algorithms execute tasks until completion without interruption. One such algorithm is the Shortest Job First (SJF) algorithm, which selects the task with the shortest processing time among all waiting tasks.
\FloatBarrier
\begin{algorithm}
\caption{Shortest Job First (SJF)}
\begin{algorithmic}[1]
\Procedure{SJF}{tasks}
\State $current\_time \gets 0$
\While{$tasks$ not empty}
\State $next\_task \gets$ task with shortest processing time
\State $execute(next\_task)$
\State $current\_time \gets current\_time + next\_task.duration$
\EndWhile
\EndProcedure
\end{algorithmic}
\end{algorithm}

\subsubsection{Comparison of Preemptive vs Non-Preemptive Scheduling}

Preemptive and non-preemptive scheduling are two contrasting approaches to task scheduling, each with its advantages and disadvantages.

\textbf{Preemptive Scheduling:}
\begin{itemize}
    \item \textbf{Responsiveness:} Preemptive scheduling offers better responsiveness as higher-priority tasks can be executed immediately when they become runnable. This is particularly beneficial in time-critical systems where tasks need to respond quickly to external events.
    \item \textbf{Dynamic Priority Adjustment:} Preemptive scheduling allows for dynamic adjustment of task priorities based on changing system conditions. Tasks with higher priority can preempt lower-priority tasks, ensuring that critical tasks are executed promptly.
    \item \textbf{Overhead:} However, preemptive scheduling introduces overhead due to frequent context switches. When a higher-priority task preempts a lower-priority task, the processor must save the state of the preempted task and restore the state of the preempting task, leading to increased overhead.
\end{itemize}

\textbf{Non-Preemptive Scheduling:}
\begin{itemize}
    \item \textbf{Predictability:} Non-preemptive scheduling offers better predictability as tasks are executed until completion without interruption. This can simplify task management and analysis, especially in real-time systems where deterministic behavior is crucial.
    \item \textbf{Reduced Overhead:} Non-preemptive scheduling reduces overhead compared to preemptive scheduling since there are no context switches during task execution. This can lead to improved system efficiency, especially on systems with limited resources.
    \item \textbf{Response Time Variability:} However, non-preemptive scheduling may result in longer response times for high-priority tasks if they are blocked by lower-priority tasks. Tasks with higher priority must wait for lower-priority tasks to complete before they can be executed, leading to increased variability in response times.
\end{itemize}

\textbf{Choosing Between Preemptive and Non-Preemptive Scheduling:}
The choice between preemptive and non-preemptive scheduling depends on the specific requirements of the system and the characteristics of the tasks being scheduled. In time-critical systems where responsiveness is paramount, preemptive scheduling may be preferred despite the overhead. In contrast, for systems where predictability and determinism are more important, non-preemptive scheduling may be a better choice to ensure consistent behavior.

In practice, hybrid approaches that combine preemptive and non-preemptive scheduling techniques may be used to balance responsiveness and predictability, depending on the system's requirements and constraints.


\subsection{Multiprocessor Scheduling}

Multiprocessor scheduling involves efficiently allocating tasks to multiple processors for parallel execution, aiming to maximize throughput and minimize resource contention.

\subsubsection{Multiprocessor Scheduling Algorithm}

One approach to multiprocessor scheduling is the Work-Conserving algorithm, which ensures that each processor remains busy as much as possible by assigning tasks dynamically based on their computational requirements.
\FloatBarrier
\begin{algorithm}
\caption{Work-Conserving Multiprocessor Scheduling}
\begin{algorithmic}[1]
\Procedure{Work-Conserving}{tasks, processors}
\While{$tasks$ not empty}
\State $task \gets$ select task with highest priority
\State $processor \gets$ select idle processor
\State $assign(task, processor)$
\EndWhile
\EndProcedure
\end{algorithmic}
\end{algorithm}

\subsection{Real-Time Scheduling}

Real-time scheduling is crucial for systems where tasks must meet strict deadlines to ensure proper operation, such as in embedded systems and multimedia applications.

\subsubsection{Real-Time Scheduling Algorithm}

The Rate-Monotonic Scheduling (RMS) algorithm assigns priorities to tasks based on their periods, with shorter periods assigned higher priorities. This ensures that tasks with tighter deadlines are executed before those with looser deadlines.
\FloatBarrier
\begin{algorithm}
\caption{Rate-Monotonic Scheduling (RMS)}
\begin{algorithmic}[1]
\Procedure{RMS}{tasks}
\State $current\_time \gets 0$
\While{$tasks$ not empty}
\State $next\_task \gets$ task with highest priority
\State $execute(next\_task)$
\State $current\_time \gets current\_time + next\_task.period$
\EndWhile
\EndProcedure
\end{algorithmic}
\end{algorithm}

\section{Scheduling in Operating Systems}
Operating systems employ various scheduling techniques to manage and optimize the allocation of resources such as CPU time, memory, and I/O devices. Process scheduling involves determining the order in which processes are executed by the CPU. Thread scheduling focuses on managing the execution of individual threads within a process. Additionally, I/O scheduling aims to optimize the order in which I/O requests are serviced by devices such as disks and network interfaces.

\subsection{Process Scheduling}
Process scheduling is a fundamental aspect of operating system design, as it directly impacts system performance and responsiveness. The goal of process scheduling is to maximize system throughput, minimize response time, and ensure fair allocation of resources among competing processes.

One of the most commonly used process scheduling algorithms is the \textbf{Round Robin (RR)} scheduling algorithm. RR is a preemptive scheduling algorithm that assigns a fixed time quantum to each process in the ready queue. When a process's time quantum expires, it is preempted and moved to the end of the queue, allowing the next process to execute.

Let $T_i$ represent the total execution time of process $i$, and $t_i$ represent the time quantum assigned to process $i$ in RR scheduling. The average waiting time ($W_{\text{avg}}$) for all processes in the queue can be calculated using the following formula:

\[
W_{\text{avg}} = \frac{\sum_{i=1}^{n} (W_i)}{n}
\]

Where $W_i$ represents the waiting time for process $i$, which can be calculated as:

\[
W_i = (n - 1) \times t_i
\]

The RR scheduling algorithm ensures fairness by providing each process with an equal share of CPU time, regardless of its priority or execution characteristics.
\FloatBarrier
\begin{algorithm}
\caption{Round Robin (RR) Scheduling Algorithm}\label{rr}
\begin{algorithmic}[1]
\Procedure{RoundRobin}{$processes, quantum$}
\State $queue \gets empty$
\State $time \gets 0$
\While{$\text{not } processes.empty()$}
\State $process \gets processes.pop()$
\State $execute(process, quantum)$
\State $time \gets time + quantum$
\If{$process \text{ not finished}$}
\State $queue.push(process)$
\EndIf
\EndWhile
\EndProcedure
\end{algorithmic}
\end{algorithm}
\FloatBarrier
\textbf{Python Code Implementation:}
\FloatBarrier
\begin{lstlisting}
def round_robin(processes, quantum):
    queue = []
    time = 0
    while processes:
        process = processes.pop(0)
        execute(process, quantum)
        time += quantum
        if not process.finished:
            queue.append(process)
    return
\end{lstlisting}
\FloatBarrier

\subsection{Thread Scheduling}
Thread scheduling focuses on managing the execution of individual threads within a process. One popular thread scheduling algorithm is the \textbf{Multilevel Feedback Queue (MLFQ)} algorithm. MLFQ maintains multiple queues with different priority levels, where threads are scheduled based on their priority and recent execution history.
\FloatBarrier
\begin{algorithm}
\caption{Multilevel Feedback Queue (MLFQ) Scheduling Algorithm}\label{mlfq}
\begin{algorithmic}[1]
\Procedure{MLFQ}{$threads$}
\State $queues \gets initialize\_queues()$
\While{$\text{not all queues are empty}$}
\For{$i \gets 0$ to $n$}
\If{$queue[i] \text{ not empty}$}
\State $thread \gets queue[i].pop()$
\State $execute(thread)$
\If{$thread \text{ not finished}$}
\State $queues.adjust\_priority(thread)$
\EndIf
\EndIf
\EndFor
\EndWhile
\EndProcedure
\end{algorithmic}
\end{algorithm}
\FloatBarrier
\textbf{Python Code Implementation:}
\FloatBarrier
\begin{lstlisting}
def MLFQ(threads):
    queues = initialize_queues()
    while not all_empty(queues):
        for queue in queues:
            if queue:
                thread = queue.pop(0)
                execute(thread)
                if not thread.finished:
                    queues.adjust_priority(thread)
    return
\end{lstlisting}

\subsection{I/O Scheduling}
I/O scheduling is concerned with optimizing the order in which I/O requests are serviced by devices such as disks and network interfaces. One commonly used I/O scheduling algorithm is the \textbf{Shortest Seek Time First (SSTF)} algorithm. SSTF selects the I/O request that is closest to the current disk head position, minimizing the seek time required to access data.
\FloatBarrier
\begin{algorithm}
\caption{Shortest Seek Time First (SSTF) Scheduling Algorithm}\label{sstf}
\begin{algorithmic}[1]
\Procedure{SSTF}{$requests, head\_position$}
\State $sorted\_requests \gets sort\_by\_distance(requests, head\_position)$
\For{$request$ in $sorted\_requests$}
\State $execute(request)$
\EndFor
\EndProcedure
\end{algorithmic}
\end{algorithm}
\FloatBarrier
\textbf{Python Code Implementation:}
\FloatBarrier
\begin{lstlisting}
def SSTF(requests, head_position):
    sorted_requests = sort_by_distance(requests, head_position)
    for request in sorted_requests:
        execute(request)
    return
\end{lstlisting}


\section{Scheduling in Production and Manufacturing}
Scheduling is key in production and manufacturing, determining task sequences on machines to enhance efficiency, reduce production time, and optimize resources. Here, we explore Job Shop Scheduling, Flow Shop Scheduling, and Project Scheduling.

\subsection{Job Shop Scheduling}
Job Shop Scheduling schedules jobs with multiple operations on specific machines in a set order, common in settings like automobile assembly lines and semiconductor plants.

A popular algorithm for Job Shop Scheduling is \textbf{Johnson's Rule}, which minimizes total completion time by sequencing operations optimally. It identifies the critical path, the sequence that determines overall completion time, and schedules accordingly.


\textbf{Algorithm: Johnson's Rule}
\begin{enumerate}
    \item For each job, compute the total processing time as the sum of processing times for all operations.
    \item Initialize two queues: one for machines with the smallest processing time first (Queue 1) and the other for machines with the largest processing time first (Queue 2).
    \item While both queues are not empty:
    \begin{itemize}
        \item Choose the operation with the smallest processing time from Queue 1 and schedule it next.
        \item Choose the operation with the smallest processing time from Queue 2 and schedule it next.
    \end{itemize}
\end{enumerate}

Let's illustrate Johnson's Rule with an example:
\FloatBarrier
\FloatBarrier
\begin{algorithm}
\caption{Johnson's Rule}
\begin{algorithmic}[1]
\Procedure{Johnson}{Jobs, Machines}
\State Compute total processing time for each job
\State Initialize Queues 1 and 2
\While{Both Queues are not empty}
    \State Schedule operation with smallest processing time from Queue 1
    \State Schedule operation with smallest processing time from Queue 2
\EndWhile
\EndProcedure
\end{algorithmic}
\end{algorithm}
\FloatBarrier

\subsection{Flow Shop Scheduling}
Flow Shop Scheduling involves scheduling a set of jobs that need to be processed on a series of machines in the same order, where each job consists of multiple operations and each operation must be processed on a specific machine in a predefined order. This problem arises in manufacturing settings where the sequence of operations is fixed, such as in assembly lines.

One of the most common algorithms used for Flow Shop Scheduling is the \textbf{Johnson's Algorithm}. This algorithm extends Johnson's Rule to the flow shop setting by scheduling jobs on two machines in a way that minimizes the total completion time.

\textbf{Algorithm: Johnson's Algorithm}
\begin{enumerate}
    \item For each job, compute the total processing time as the sum of processing times for all operations.
    \item Initialize two queues: one for machines with the smallest processing time first (Queue 1) and the other for machines with the largest processing time first (Queue 2).
    \item While both queues are not empty:
    \begin{itemize}
        \item Choose the operation with the smallest processing time from Queue 1 and schedule it next.
        \item Choose the operation with the smallest processing time from Queue 2 and schedule it next.
    \end{itemize}
\end{enumerate}

\subsection{Project Scheduling}
Project Scheduling involves scheduling a set of activities or tasks that need to be completed within a specified time frame, taking into account dependencies between tasks and resource constraints. This problem arises in project management, construction projects, and software development.

One of the most common algorithms used for Project Scheduling is the \textbf{Critical Path Method (CPM)}. CPM aims to determine the longest path of dependent activities, known as the critical path, and the total duration of the project. By identifying the critical path, project managers can allocate resources efficiently and ensure timely project completion.

\textbf{Algorithm: Critical Path Method (CPM)}
\begin{enumerate}
    \item Construct a network diagram representing all tasks and their dependencies.
    \item Compute the earliest start time (ES) and earliest finish time (EF) for each task.
    \item Compute the latest start time (LS) and latest finish time (LF) for each task.
    \item Calculate the slack time for each task (slack = LF - EF).
    \item Identify the critical path, which consists of tasks with zero slack time.
    \item Determine the total duration of the project as the sum of durations along the critical path.
\end{enumerate}

\section{Scheduling in Distributed Systems}

Scheduling in distributed systems allocates tasks to resources efficiently, optimizing performance and resource use. This section covers load balancing, task allocation, and cloud computing scheduling.

\subsection{Load Balancing}

Load balancing evenly distributes workloads across resources to optimize utilization, prevent overload, and enhance performance. It involves distributing tasks or requests among resources to minimize response time, maximize throughput, and ensure fairness.

\subsubsection{Need for Load Balancing}

In distributed systems, workloads are dynamic and unpredictable, with varying demand and resource availability. Effective load balancing prevents some resources from being overwhelmed while others are underutilized, thus maintaining a balanced workload across the system.


\subsubsection{Least Loaded Algorithm}

The Least Loaded Algorithm assigns incoming tasks to the resource with the least load, where load can be measured in terms of CPU utilization, memory usage, or other relevant metrics. Mathematically, the algorithm can be described as follows:

Let $R$ be the set of resources, and $L(r)$ be the load on resource $r \in R$. When a new task $T$ arrives, the algorithm selects the resource $r^*$ such that:

$$r^* = \text{argmin}_{r \in R} L(r)$$

The task $T$ is then assigned to the resource $r^*$.
\FloatBarrier
\begin{algorithm}[H]
\caption{Least Loaded Algorithm}\label{alg:least_loaded}
\begin{algorithmic}[1]
\Procedure{LeastLoaded}{$R, T$}
\State $r^* \gets \text{argmin}_{r \in R} L(r)$
\State Assign task $T$ to resource $r^*$
\EndProcedure
\end{algorithmic}
\end{algorithm}
\FloatBarrier   
\textbf{Python Code Implementation:}
\FloatBarrier
\begin{lstlisting}
def least_loaded(resources, task):
    least_loaded_resource = min(resources, key=lambda r: r.load)
    least_loaded_resource.assign_task(task)
\end{lstlisting}
\FloatBarrier
The Least Loaded Algorithm selects the resource with the minimum load and assigns the incoming task to that resource. This helps in achieving load balancing across distributed systems efficiently.

\subsection{Task Allocation}

Task allocation is a critical aspect of distributed systems where tasks need to be assigned to resources in a way that optimizes system performance, resource utilization, and overall efficiency. In this subsection, we will delve deeper into the concept of task allocation and explore additional task allocation algorithms.

Task allocation algorithms aim to assign tasks to resources based on various factors such as task characteristics, resource capabilities, and system constraints. These algorithms help optimize resource utilization, minimize task completion time, and ensure fairness in task distribution.

\subsubsection{Need for Task Allocation}

In distributed systems, tasks may have different requirements and constraints, and resources may vary in terms of capacity, performance, and availability. Effective task allocation is essential to ensure that tasks are executed on appropriate resources to meet performance objectives, resource constraints, and service-level agreements.

\subsubsection{Task Allocation Algorithm}

One common approach to task allocation is the Max-Min Algorithm, which aims to maximize the minimum resource allocation for each task.
\FloatBarrier
\begin{algorithm}[H]
\caption{Max-Min Algorithm}\label{alg:max_min}
\begin{algorithmic}[1]
\Procedure{MaxMin}{$T, R$}
\State Sort tasks in descending order of resource requirements
\State Sort resources in ascending order of available capacity
\For{each task $t$}
    \State Allocate $t$ to the resource with the most available capacity
\EndFor
\EndProcedure
\end{algorithmic}
\end{algorithm}
\FloatBarrier   
\textbf{Python Code Implementation:}
\FloatBarrier
\begin{lstlisting}
def max_min(tasks, resources):
    tasks.sort(key=lambda t: t.requirements, reverse=True)
    resources.sort(key=lambda r: r.available_capacity)
    
    for task in tasks:
        resource = max(resources, key=lambda r: r.available_capacity)
        resource.assign_task(task)
\end{lstlisting}
\FloatBarrier
The Max-Min Algorithm optimizes task allocation by maximizing the minimum resource allocation for each task, thereby improving system performance and resource utilization.

\subsection{Cloud Computing Scheduling}

Cloud computing scheduling involves efficiently allocating tasks to virtual machines (VMs) or containers in cloud environments to meet service-level objectives and minimize costs. In this subsection, we will delve deeper into the concept of cloud computing scheduling and explore additional scheduling algorithms.

\subsubsection{Need for Cloud Computing Scheduling}

Cloud computing environments typically host a large number of VMs or containers, each running various tasks or services. Efficiently allocating tasks to these resources is essential to maximize resource utilization, ensure performance scalability, and optimize cost-effectiveness. Cloud computing scheduling helps address these challenges by dynamically assigning tasks to resources based on workload characteristics, resource availability, and system requirements.


\subsubsection{Cloud Computing Scheduling Algorithm}

One widely used scheduling algorithm in cloud computing is the First Fit Decreasing (FFD) algorithm, which sorts tasks in descending order of resource requirements and allocates them to the first VM or container that has sufficient capacity.
\FloatBarrier
\begin{algorithm}[H]
\caption{First Fit Decreasing (FFD) Algorithm}\label{alg:ffd}
\begin{algorithmic}[1]
\Procedure{FFD}{$T, VMs$}
\State Sort tasks in descending order of resource requirements
\For{each task $t$}
    \For{each VM $v$}
        \If{$t$ fits in $v$}
            \State Allocate $t$ to $v$
            \State Break
        \EndIf
    \EndFor
\EndFor
\EndProcedure
\end{algorithmic}
\end{algorithm}
\FloatBarrier
\textbf{Python Code Implementation:}
\FloatBarrier
\begin{lstlisting}
def ffd(tasks, vms):
    tasks.sort(key=lambda t: t.requirements, reverse=True)
    
    for task in tasks:
        for vm in vms:
            if task.fits(vm):
                vm.assign_task(task)
                break
\end{lstlisting}
\FloatBarrier
The FFD algorithm efficiently allocates tasks to VMs in cloud environments, helping optimize resource utilization and meet service-level objectives effectively.


\section{Scheduling in Networks}

Scheduling in networks plays a critical role in optimizing the allocation of resources and ensuring efficient communication. This section explores various aspects of scheduling techniques in network systems, including packet scheduling in network routers, bandwidth allocation, and quality of service (QoS) management.

\subsection{Packet Scheduling in Network Routers}

Packet scheduling in network routers is a fundamental aspect of traffic management, where incoming packets are scheduled for transmission based on certain criteria. This ensures fair distribution of resources and minimizes congestion in the network.

Packet scheduling algorithms aim to allocate resources efficiently while meeting certain performance criteria such as delay, throughput, and fairness. One widely used packet scheduling algorithm is Weighted Fair Queuing (WFQ).

\subsubsection{Mathematical Discussion}

Let $N$ be the number of flows in the network, and let $w_i$ denote the weight assigned to flow $i$, where $1 \leq i \leq N$. The weight represents the relative priority or importance of each flow. Let $R_i$ be the rate allocated to flow $i$, representing the amount of bandwidth assigned to the flow.

The WFQ algorithm assigns a virtual finish time $F_i$ to each packet in flow $i$, calculated as:

\begin{equation*}
    F_i = \text{Arrival time} + \frac{\text{Packet size}}{R_i}
\end{equation*}

The packets are then enqueued in a priority queue based on their virtual finish times, and transmitted in the order of increasing virtual finish times.

\subsubsection{Algorithmic Example: Weighted Fair Queuing (WFQ)}
\FloatBarrier
\begin{algorithm}[H]
\caption{Weighted Fair Queuing (WFQ)}
\begin{algorithmic}[1]
\State Initialize virtual time for each flow
\State Initialize empty priority queue
\While{packets arrive}
    \State Assign virtual finish time for each packet based on its arrival time and service rate
    \State Enqueue packets into priority queue based on virtual finish time
    \State Dequeue and transmit packets in order of increasing virtual finish time
\EndWhile
\end{algorithmic}
\end{algorithm}
\FloatBarrier

\subsection{Bandwidth Allocation}

Bandwidth allocation involves dividing the available bandwidth among competing users or applications in a network. This ensures that each user receives a fair share of the available resources while maximizing overall network efficiency.

\subsubsection{Mathematical Discussion}

In bandwidth allocation, the goal is to distribute the available bandwidth among $N$ users or applications in a fair and efficient manner. Let $X_i$ denote the bandwidth allocated to user $i$, where $1 \leq i \leq N$. The total available bandwidth is denoted by $B$.

One common approach to bandwidth allocation is Max-Min Fairness. In this approach, the bandwidth allocation vector $\mathbf{X}$ is iteratively adjusted until convergence. At each iteration, the deficit for each user is calculated based on the current allocation. The deficit represents the difference between the user's allocated bandwidth and its fair share. Bandwidth is then allocated to each user proportional to its deficit until convergence is achieved.

\subsubsection{Algorithmic Example: Max-Min Fairness}
\FloatBarrier
\begin{algorithm}[H]
\caption{Max-Min Fairness}
\begin{algorithmic}[1]
\State Initialize bandwidth allocation vector $\mathbf{X}$ with equal shares
\While{not converged}
    \State Calculate the deficit for each user based on current allocation
    \State Allocate bandwidth to each user proportional to its deficit until convergence
\EndWhile
\end{algorithmic}
\end{algorithm}
\FloatBarrier

\subsection{Quality of Service (QoS) Management}

Quality of Service (QoS) management involves prioritizing certain types of traffic or users to meet specific performance requirements, such as latency, throughput, and reliability. This ensures that critical applications receive the necessary resources to function optimally.

\subsubsection{Mathematical Discussion}

In QoS management, different traffic classes are assigned different priority levels based on their QoS requirements. Let $T_i$ denote the traffic class of flow $i$, where $1 \leq i \leq N$. Each traffic class is associated with a set of performance parameters, such as latency bounds, packet loss rates, and throughput requirements.

One approach to QoS management is the Token Bucket Algorithm. In this algorithm, a token bucket is used to control the rate at which packets are transmitted. Each packet is allowed to be transmitted only if there are sufficient tokens in the bucket. If tokens are depleted, packets are either dropped or queued for later transmission.

\subsubsection{Algorithmic Example: Token Bucket Algorithm}
\FloatBarrier
\begin{algorithm}[H]
\caption{Token Bucket Algorithm}
\begin{algorithmic}[1]
\State Initialize token bucket with tokens and token replenishment rate
\While{packets arrive}
    \If{tokens available}
        \State Transmit packet
        \State Consume tokens
    \Else
        \State Drop packet
    \EndIf
\EndWhile
\end{algorithmic}
\end{algorithm}
\FloatBarrier

\section{Mathematical and Heuristic Methods in Scheduling}

Scheduling problems involve allocating limited resources over time to perform a set of tasks, subject to various constraints and objectives. Mathematical and heuristic methods are commonly used to solve scheduling problems efficiently and effectively. In this section, we explore different approaches, including Linear Programming, Heuristic Algorithms, and Evolutionary Algorithms.

\subsection{Linear Programming Approaches}

Linear Programming (LP) provides a mathematical framework for optimizing linear objective functions subject to linear constraints. In the context of scheduling, LP approaches formulate the scheduling problem as a linear optimization problem and solve it using LP techniques.

\begin{itemize}
    \item \textbf{Job Scheduling with Minimization of Total Completion Time:} In this approach, the objective is to minimize the total completion time of all jobs subject to various constraints such as job dependencies and resource availability. Mathematically, this can be formulated as follows:
    \FloatBarrier
    \begin{align*}
        & \text{Minimize} \quad \sum_{i=1}^{n} C_i \\
        & \text{Subject to} \\
        & C_i = \text{start time of job } i \\
        & C_i + p_i \leq C_j \quad \forall (i, j) \in \text{precedence constraints} \\
        & C_i + p_i \leq d_i \quad \forall i \in \text{jobs} \\
        & C_i \geq 0 \quad \forall i \in \text{jobs}
    \end{align*}
    \FloatBarrier
    where \(C_i\) represents the completion time of job \(i\), \(p_i\) is the processing time of job \(i\), and \(d_i\) is the due date of job \(i\). The first constraint ensures that each job starts after its predecessor completes, the second constraint enforces the due date constraint, and the last constraint ensures non-negative start times.
    
    \item \textbf{Job Scheduling with Minimization of Total Weighted Tardiness:} This approach aims to minimize the total weighted tardiness of jobs, where tardiness is the amount of time a job finishes after its due date, weighted by a given weight for each job. Mathematically, this can be formulated as follows:
    \FloatBarrier
    \begin{align*}
        & \text{Minimize} \quad \sum_{i=1}^{n} w_i \cdot \max(0, C_i - d_i) \\
        & \text{Subject to} \\
        & C_i = \text{start time of job } i \\
        & C_i + p_i \leq C_j \quad \forall (i, j) \in \text{precedence constraints} \\
        & C_i + p_i \leq d_i \quad \forall i \in \text{jobs} \\
        & C_i \geq 0 \quad \forall i \in \text{jobs}
    \end{align*}
    \FloatBarrier
    where \(w_i\) is the weight of job \(i\), and all other constraints are similar to the previous formulation.
\end{itemize}

Let's illustrate the Job Scheduling with Minimization of Total Completion Time using the following algorithmic example:
\FloatBarrier
\begin{algorithm}
\caption{LP-based Job Scheduling with Minimization of Total Completion Time}
\begin{algorithmic}[1]
\State Formulate the LP problem as described above.
\State Solve the LP problem using an LP solver (e.g., simplex method).
\State Extract the optimal solution to obtain the start times of each job.
\end{algorithmic}
\end{algorithm}
\FloatBarrier

\subsection{Heuristic Algorithms}

Heuristic algorithms provide approximate solutions to scheduling problems in a reasonable amount of time. These algorithms do not guarantee optimality but often produce high-quality solutions efficiently.

\begin{itemize}
    \item \textbf{Earliest Due Date (EDD) Algorithm:} The EDD algorithm schedules jobs based on their due dates, with the earliest due date job scheduled first. Mathematically, the EDD rule can be expressed as:
    \[ \text{Priority}(i) = d_i \]
    where \(d_i\) is the due date of job \(i\).
    
    \item \textbf{Shortest Processing Time (SPT) Algorithm:} The SPT algorithm schedules jobs based on their processing times, with the shortest processing time job scheduled first. Mathematically, the SPT rule can be expressed as:
    \[ \text{Priority}(i) = p_i \]
    where \(p_i\) is the processing time of job \(i\).
\end{itemize}

Let's demonstrate the Earliest Due Date (EDD) Algorithm using the following algorithmic example:
\FloatBarrier
\begin{algorithm}
\caption{Earliest Due Date (EDD) Algorithm}
\begin{algorithmic}[1]
\For{each job \(i\) in non-decreasing order of due dates}
\State Schedule job \(i\) at the earliest available time slot.
\EndFor
\end{algorithmic}
\end{algorithm}
\FloatBarrier

\subsection{Evolutionary Algorithms and Machine Learning}

Evolutionary algorithms (EAs) are population-based optimization techniques inspired by the process of natural selection. These algorithms iteratively evolve a population of candidate solutions to search for optimal or near-optimal solutions to a given problem. In the context of scheduling, EAs are used to explore the solution space efficiently and find high-quality schedules.

\begin{itemize}
    \item \textbf{Genetic Algorithm (GA):} GA is one of the most well-known EAs. It operates on a population of potential solutions, where each solution is represented as a chromosome. The chromosome encodes a candidate solution to the scheduling problem, typically representing a permutation of job orders or a binary encoding of job allocations. GA iteratively applies genetic operators such as selection, crossover, and mutation to evolve better solutions over generations.

    \item \textbf{Particle Swarm Optimization (PSO):} PSO is another popular EA based on the social behavior of particles in a swarm. In PSO, each potential solution is represented as a particle moving in the search space. The position of each particle corresponds to a candidate solution, and the velocity of each particle determines how it explores the search space. PSO iteratively updates the position and velocity of particles based on their own best-known position and the global best-known position in the search space.

\end{itemize}

Let's provide more detailed algorithmic examples for both GA and PSO in the context of scheduling:

\subsubsection{Genetic Algorithm (GA)}

In the context of scheduling, a chromosome in GA represents a permutation of job orders. The objective is to find the best permutation that minimizes a scheduling criterion such as total completion time or total weighted tardiness.
\FloatBarrier
\begin{algorithm}
\caption{Genetic Algorithm for Job Scheduling}
\begin{algorithmic}[1]
\State Initialize population of chromosomes randomly.
\For{each generation}
    \State Evaluate fitness of each chromosome.
    \State Select parents for reproduction based on fitness.
    \State Apply crossover and mutation operators to create offspring.
    \State Replace old population with new population of offspring.
\EndFor
\Return Best chromosome found.
\end{algorithmic}
\end{algorithm}

\subsubsection{Particle Swarm Optimization (PSO)}

In PSO, each particle represents a potential solution, and its position corresponds to a candidate schedule. The objective is to find the best position (schedule) that minimizes a scheduling criterion.
\FloatBarrier
\begin{algorithm}
\caption{Particle Swarm Optimization for Job Scheduling}
\begin{algorithmic}[1]
\State Initialize particles with random positions and velocities.
\For{each iteration}
    \State Update personal best-known position for each particle.
    \State Update global best-known position for the swarm.
    \State Update particle positions and velocities.
\EndFor
\Return Best position found.
\end{algorithmic}
\end{algorithm}
\FloatBarrier
Both GA and PSO can be customized and extended to handle various scheduling criteria and constraints. They offer the advantage of exploring the solution space effectively and finding high-quality schedules, although they may not guarantee optimality.

Evolutionary algorithms are relevant to machine learning in scheduling as they can adaptively learn from past schedules and improve scheduling policies over time, leading to more efficient and effective scheduling solutions.


\section{Challenges and Future Directions}

Scheduling is essential across various domains, but as systems grow more complex, new challenges and opportunities arise. This section discusses key challenges and future directions in scheduling research.

\subsection{Scalability Issues}

Scalability is a critical concern in scheduling algorithms, especially as systems grow larger and more complex. The computational complexity of scheduling problems often grows exponentially with the size of the problem. 

Consider a scheduling problem with \(n\) tasks and \(m\) resources. Let \(T = \{1, 2, ..., n\}\) be the set of tasks and \(R = \{1, 2, ..., m\}\) the set of resources. Each task \(i \in T\) requires \(r_{ij}\) of each resource \(j \in R\). We introduce binary decision variables \(x_{ijt}\), where \(x_{ijt} = 1\) if task \(i\) is assigned to resource \(j\) at time \(t\), and \(x_{ijt} = 0\) otherwise. 

The goal is to minimize or maximize an objective function subject to constraints. For example, to minimize makespan:
\[
\text{Minimize } f(x_{ijt})
\]
subject to:
\[
\sum_{j \in R} x_{ijt} \leq C_{ij} \quad \forall i \in T, \forall t
\]

Many scheduling problems are NP-hard, making them computationally intractable for large instances. Heuristic or approximation algorithms, such as greedy algorithms and metaheuristic techniques, provide near-optimal solutions efficiently.

\subsection{Scheduling under Uncertainty}

Scheduling under uncertainty involves making decisions with incomplete information, such as variable task durations and resource availability. This can be modeled using stochastic optimization techniques, probabilistic models, and decision theory.

For example, consider a scheduling problem with uncertain task durations modeled as stochastic optimization. Using dynamic programming, we can find a schedule that minimizes expected completion time or maximizes resource utilization while accounting for uncertainty.

\subsection{Emerging Applications and Technologies}

Scheduling techniques are being applied to new domains and technologies, driving research and development:

\begin{itemize}
    \item \textbf{Smart Manufacturing}: Optimizing production processes in interconnected systems to improve efficiency and reduce downtime.
    \item \textbf{Internet of Things (IoT)}: Coordinating resources and tasks in IoT networks with limited capabilities.
    \item \textbf{Blockchain Technology}: Scheduling transactions and smart contracts in decentralized networks.
    \item \textbf{Edge Computing}: Optimizing task scheduling and resource allocation closer to end-users to reduce latency.
    \item \textbf{Autonomous Systems}: Planning and scheduling tasks for autonomous vehicles, drones, and robots in dynamic environments.
\end{itemize}

These applications present unique challenges and opportunities, driving the development of new algorithms and optimization methods tailored to specific domains and technologies.


\section{Case Studies}
Scheduling techniques are essential in various domains, including large-scale data centers, smart grids, and autonomous systems. This section explores scheduling challenges, algorithms, and applications in these contexts.

\subsection{Scheduling in Large-Scale Data Centers}
Large-scale data centers host numerous applications and services, making efficient scheduling crucial for optimizing resource utilization, minimizing latency, and maximizing throughput. 

\subsubsection{Mathematically Detailed Discussion}
In data centers, tasks arrive dynamically and must be assigned to available resources such as servers or virtual machines (VMs). The scheduling problem can be formulated as an optimization task to minimize total completion time or maximize resource utilization while meeting constraints like deadlines or service-level agreements.

\[
\text{Minimize} \sum_{i=1}^{n} C_i
\]
subject to:
\[
C_i \leq D_i \quad \forall i \in \{1, 2, \ldots, n\}
\]
where \(C_i\) is the completion time of task \(i\) and \(D_i\) is its deadline.


One commonly used algorithm for scheduling tasks in large-scale data centers is the \textbf{Minimum Completion Time (MCT)} algorithm. The MCT algorithm assigns each task to the resource that minimizes its completion time, considering factors such as processing power, network bandwidth, and data locality. The algorithm can be described as follows:
\FloatBarrier
\begin{algorithm}[H]
\caption{Minimum Completion Time (MCT) Algorithm}
\begin{algorithmic}[1]
\State Sort tasks in non-decreasing order of processing time
\For{each task $T$}
\State Assign $T$ to the resource that minimizes its completion time
\EndFor
\end{algorithmic}
\end{algorithm}

\subsubsection{Algorithmic Example}
Here's a Python implementation of the MCT algorithm for scheduling tasks in a large-scale data center:
\FloatBarrier
\begin{lstlisting}
def mct_algorithm(tasks, resources):
    tasks.sort(key=lambda x: x.processing_time)
    for task in tasks:
        min_completion_time = float('inf')
        selected_resource = None
        for resource in resources:
            completion_time = calculate_completion_time(task, resource)
            if completion_time < min_completion_time:
                min_completion_time = completion_time
                selected_resource = resource
        assign_task_to_resource(task, selected_resource)
\end{lstlisting}
\FloatBarrier

\subsection{Dynamic Scheduling in Smart Grids}
Smart grids use advanced communication and control technologies to manage electricity generation, distribution, and consumption efficiently. Dynamic scheduling optimizes energy usage, reduces costs, and maintains grid stability.

\subsubsection{Mathematically Detailed Discussion}
Dynamic scheduling in smart grids involves real-time allocation of resources like power generators, storage systems, and demand-side resources to meet electricity demand while adhering to operational constraints and market dynamics.

\[
\text{Minimize} \sum_{t=1}^{T} (C_t \cdot E_t + P_t)
\]
subject to:
\[
\sum_{i=1}^{n} G_i(t) \geq D(t) \quad \forall t
\]
where \(C_t\) is the cost of energy at time \(t\), \(E_t\) is the energy consumption, \(P_t\) is the penalty for not meeting demand, \(G_i(t)\) is the power generated by source \(i\) at time \(t\), and \(D(t)\) is the demand at time \(t\).

\textbf{Model Predictive Control (MPC)} is a technique used for dynamic scheduling in smart grids. MPC formulates the scheduling problem as a dynamic optimization task and solves it iteratively based on predictions of future behavior, adjusting resource allocations in response to changing conditions like fluctuating demand or renewable energy availability.


\subsubsection{Algorithmic Example}
Here's a simplified Python implementation of the Model Predictive Control (MPC) algorithm for dynamic scheduling in a smart grid:
\FloatBarrier
\begin{lstlisting}
def mpc_algorithm():
    initialize_system_state()
    for t in range(total_time_steps):
        predict_future_state()
        optimize_resource_allocation()
        update_system_state()
\end{lstlisting}
\FloatBarrier

\subsection{Adaptive Scheduling in Autonomous Systems}
Autonomous systems like unmanned aerial vehicles (UAVs) and self-driving cars use adaptive scheduling to make real-time decisions and coordinate actions efficiently.

\subsubsection{Mathematically Detailed Discussion}
Adaptive scheduling in autonomous systems involves dynamically allocating tasks or resources based on changing environmental conditions, system states, and mission objectives. The problem often involves uncertainty and requires robust algorithms capable of adapting to unforeseen events.

\textbf{Reinforcement Learning (RL)} is a common approach for adaptive scheduling in autonomous systems. RL formulates the scheduling problem as a Markov Decision Process (MDP) and learns optimal scheduling policies through trial and error interactions with the environment. The algorithm explores different strategies and updates its policy based on observed rewards and penalties.

\[
\pi^* = \arg\max_\pi \mathbb{E}\left[\sum_{t=0}^{T} \gamma^t R(s_t, a_t) \mid \pi\right]
\]
where \(\pi^*\) is the optimal policy, \(R(s_t, a_t)\) is the reward received at time \(t\) for taking action \(a_t\) in state \(s_t\), and \(\gamma\) is the discount factor.


\subsubsection{Algorithmic Example}
Here's a simple Python implementation of a Reinforcement Learning (RL) algorithm for adaptive scheduling in an autonomous system:
\FloatBarrier
\begin{lstlisting}
def reinforcement_learning():
    initialize_q_table()
    for episode in range(total_episodes):
        observe_environment()
        while not is_terminal_state():
            choose_action_based_on_policy()
            take_action()
            update_q_table()
\end{lstlisting}
\FloatBarrier

\section{Conclusion}
We explored various scheduling techniques, highlighting their significance in algorithms and potential future research directions.

\subsection{The Evolution of Scheduling Theory and Practice}
Scheduling theory has significantly evolved, driven by advances in computer science, optimization, and operations research. Initially focused on job shop scheduling, it now spans diverse domains such as project scheduling, cloud computing, and manufacturing. Early heuristic algorithms have given way to sophisticated methods using mathematical programming and combinatorial optimization. Today's advancements in AI, machine learning, and IoT are driving the development of adaptive and predictive scheduling algorithms that dynamically optimize decisions in real-time.

\subsection{Implications for Future Research and Applications}
The future of scheduling research is promising, with opportunities to advance the field and address real-world challenges:

\begin{itemize}
\item \textbf{Intelligent Scheduling Algorithms:} Future research will focus on AI and machine learning to develop algorithms that enhance decision-making and optimize scheduling objectives.
\item \textbf{Uncertainty and Risk Management:} Incorporating mechanisms to handle uncertainty and risk will enable robust scheduling in dynamic environments.
\item \textbf{Scalability and Efficiency:} There's a need for scalable algorithms that can manage large problem instances and optimize in real-time.
\end{itemize}

Scheduling algorithms have broad applications across various domains:

\begin{itemize}
\item \textbf{Manufacturing:} Optimize production schedules, minimize setup times, and maximize resource utilization.
\item \textbf{Transportation and Logistics:} Optimize routes, schedule maintenance, and coordinate deliveries.
\item \textbf{Healthcare:} Optimize patient appointments, allocate resources, and manage facilities efficiently.
\end{itemize}

These examples underscore the importance of scheduling algorithms in improving resource allocation and operational efficiency across diverse fields.

\section{Exercises and Problems}
In this section, we present a variety of exercises and problems related to scheduling algorithm techniques. These exercises aim to reinforce the concepts covered in the chapter and provide opportunities for students to apply their knowledge in practical scenarios. The section is divided into two subsections: Conceptual Questions to Test Understanding and Practical Scenarios for Algorithm Design.

\subsection{Conceptual Questions to Test Understanding}
This subsection contains a series of conceptual questions designed to test the reader's understanding of scheduling algorithm techniques. These questions cover key concepts, properties, and theoretical aspects of various scheduling algorithms.

\begin{itemize}
    \item Explain the difference between preemptive and non-preemptive scheduling algorithms.
    \item Describe the criteria used to evaluate the performance of scheduling algorithms.
    \item Discuss the advantages and disadvantages of First-Come, First-Served (FCFS) scheduling.
    \item Compare and contrast Shortest Job First (SJF) and Shortest Remaining Time (SRT) scheduling algorithms.
    \item Explain how Round Robin (RR) scheduling works and discuss its time quantum parameter.
\end{itemize}

\subsection{Practical Scenarios for Algorithm Design}
In this subsection, we present practical scenarios related to scheduling algorithm techniques. These scenarios involve real-world problems that require the design and implementation of scheduling algorithms to optimize resource allocation and improve system performance.

\begin{itemize}
    \item \textbf{Scenario 1: CPU Scheduling}
    
    You are tasked with implementing a CPU scheduling algorithm for a multi-tasking operating system. Design an algorithm that minimizes average waiting time and ensures fairness among processes.
    
    \begin{algorithm}[H]
        \caption{Shortest Job First (SJF) Scheduling}
        \begin{algorithmic}[1]
            \Procedure{SJF}{}
            \State Sort processes by burst time in ascending order
            \State Initialize total\_time = 0, total\_waiting\_time = 0
            \For{each process in sorted processes}
                \State total\_time += process.burst\_time
                \State total\_waiting\_time += total\_time - process.arrival\_time - process.burst\_time
            \EndFor
            \State average\_waiting\_time = total\_waiting\_time / number\_of\_processes
            \State \textbf{return} average\_waiting\_time
            \EndProcedure
        \end{algorithmic}
    \end{algorithm}
    
    \FloatBarrier
    \begin{lstlisting}[language=Python, caption=Python Implementation of SJF Scheduling Algorithm]
    def sjf_scheduling(processes):
        processes.sort(key=lambda x: x.burst_time)
        total_time = 0
        total_waiting_time = 0
        for process in processes:
            total_time += process.burst_time
            total_waiting_time += total_time - process.arrival_time - process.burst_time
        average_waiting_time = total_waiting_time / len(processes)
        return average_waiting_time
    \end{lstlisting}
    \FloatBarrier
    
    \item \textbf{Scenario 2: Real-Time Task Scheduling}
    
    Your task is to develop a real-time task scheduling algorithm for an embedded system. The system has multiple tasks with different execution times and deadlines. Design an algorithm that meets all task deadlines while maximizing system throughput.
    
    \begin{algorithm}[H]
        \caption{Rate-Monotonic Scheduling}
        \begin{algorithmic}[1]
            \Procedure{RateMonotonic}{}
            \State Sort tasks by their periods in ascending order
            \State Assign priorities to tasks based on their periods (shorter periods have higher priorities)
            \State Execute tasks in order of priority, preempting lower-priority tasks if necessary
            \EndProcedure
        \end{algorithmic}
    \end{algorithm}
    
    \FloatBarrier
    \begin{lstlisting}[language=Python, caption=Python Implementation of Rate-Monotonic Scheduling Algorithm]
    def rate_monotonic_scheduling(tasks):
        tasks.sort(key=lambda x: x.period)
        for task in tasks:
            execute_task(task)
    \end{lstlisting}
    \FloatBarrier
    
    \item \textbf{Scenario 3: Disk Scheduling}
    
    Develop a disk scheduling algorithm for a hard disk drive to optimize disk access time. The algorithm should minimize disk head movement and ensure fair access to disk requests.
    
    \begin{algorithm}[H]
        \caption{SCAN (Elevator) Disk Scheduling}
        \begin{algorithmic}[1]
            \Procedure{SCAN}{}
            \State Sort disk requests by cylinder number
            \State Initialize head direction to right
            \State Traverse requests in sorted order, servicing requests until the end
            \State Change head direction when reaching the last request
            \State Continue traversing in the opposite direction until all requests are serviced
            \EndProcedure
        \end{algorithmic}
    \end{algorithm}
    
    \FloatBarrier
    \begin{lstlisting}[language=Python, caption=Python Implementation of SCAN Disk Scheduling Algorithm]
    def scan_disk_scheduling(requests):
        requests.sort()
        head_direction = 'right'
        while requests:
            if head_direction == 'right':
                for request in requests:
                    service_request(request)
                head_direction = 'left'
            else:
                for request in reversed(requests):
                    service_request(request)
                head_direction = 'right'
        \end{lstlisting}
    \FloatBarrier
\end{itemize}

\section{Further Reading and Resources}

When diving deeper into scheduling algorithm techniques, it's essential to explore various resources to gain a comprehensive understanding. Below are some valuable resources categorized into foundational texts and influential papers, online courses and tutorials, and software and tools for scheduling analysis.

\subsection{Foundational Texts and Influential Papers}

Understanding the foundational principles and key research papers is crucial for mastering scheduling algorithm techniques. Here's a curated list of influential texts and papers:

\begin{itemize}
    \item \textbf{Foundational Texts}:
    \begin{itemize}
        \item "Introduction to Algorithms" by Thomas H. Cormen, Charles E. Leiserson, Ronald L. Rivest, and Clifford Stein.
        \item "Scheduling: Theory, Algorithms, and Systems" by Michael Pinedo.
        \item "Operating System Concepts" by Abraham Silberschatz, Peter Baer Galvin, and Greg Gagne.
    \end{itemize}
    
    \item \textbf{Influential Papers}:
    \begin{itemize}
        \item "Optimal Control of Discrete-Event Systems" by N. N. Ulanov, M. I. Kudenko, and G. V. Tsarev.
        \item "Analysis of Multilevel Feedback Queues" by W. A. Halang and G. B. Halang.
        \item "Real-Time Scheduling Theory: A Historical Perspective" by C. L. Liu and J. W. Layland.
    \end{itemize}
\end{itemize}

\subsection{Online Courses and Tutorials}

Online courses and tutorials offer interactive learning experiences and in-depth explanations of scheduling algorithms. Here are some recommended resources:

\begin{itemize}
    \item \textbf{Coursera}:
    \begin{itemize}
        \item "Operating Systems" by University of London.
        \item "Discrete Optimization" by The University of Melbourne.
    \end{itemize}
    
    \item \textbf{edX}:
    \begin{itemize}
        \item "Algorithm Design and Analysis" by University of Pennsylvania.
        \item "Real-Time Systems" by Karlsruhe Institute of Technology.
    \end{itemize}
    
    \item \textbf{YouTube Tutorials}:
    \begin{itemize}
        \item "Introduction to Scheduling Algorithms" by MyCodeSchool.
        \item "Real-Time Operating Systems Concepts" by Gate Lectures by Ravindrababu Ravula.
    \end{itemize}
\end{itemize}

\subsection{Software and Tools for Scheduling Analysis}

Utilizing specialized software and tools can streamline scheduling analysis tasks and aid in implementing and testing scheduling algorithms. Here are some popular options:

\begin{itemize}
    \item \textbf{SimPy}: SimPy is a discrete-event simulation library for Python. It provides components for modeling and simulating various scheduling scenarios.
    
    \item \textbf{GanttProject}: GanttProject is an open-source project scheduling and management tool. It allows users to create Gantt charts to visualize and manage scheduling tasks.
    
    \item \textbf{MATLAB}: MATLAB offers robust tools for numerical computing and simulation. It includes functions and toolboxes for modeling and analyzing scheduling algorithms.
    
    \item \textbf{SimEvents}: SimEvents is a discrete-event simulation toolset in MATLAB that extends Simulink for modeling and analyzing event-driven systems, including scheduling processes.
\end{itemize}
\end{document}